% Section 5 — The bounded-learner gap (charter §4.5; issues T1-3, T2-4)
\section{The bounded-learner gap}
\label{sec:bounded}

% STUB (T1-3, T2-4). The positive half of the thesis (charter §1).
% Area role (charter §2): where format can help once information is held
% constant; an UPPER BOUND on achievable representation advantage via
% capacity/sample complexity; the regime that activates it.
%
% Plan:
%   - T1-3: statistical learning theory foundations (hypothesis class
%     \hypclass, capacity \capac{\hypclass}, Rademacher \Rad{\hypclass},
%     empirical risk \emprisk).
%   - T2-4: the bounded-learner advantage upper bound: the representation
%     advantage \advantage is controlled by a capacity/complexity term, not
%     by information content. This bounds effect SIZE and is the antidote to
%     "a math-dressed null" (guardrail G2).
%
% The positive content is delivered as falsifiable predictions in §8, never
% as established findings (guardrail G2).

\noindent\emph{[Stub. Learning-theory foundations in T1-3; advantage upper
bound in T2-4.]}

\begin{proposition}[Bounded-regime advantage is capacity-controlled --- placeholder]
\label{prop:advantage}
\emph{[Placeholder for T2-4.]} The achievable representation advantage
$\advantage$ at finite capacity satisfies an upper bound in terms of
$\capac{\hypclass}$ and the sample size $\nsample$.
\idealization{Holds in the bounded regime (finite capacity / finite sample
/ bounded compute); it does not assert any effect exists, only bounds its
maximum magnitude.}
\end{proposition}
